\section{Reproducibility Details}
\label{app:repro}

\subsection{Compute}
\label{app:compute}

Experiments use two GPUs depending on model scale:
\begin{itemize}[leftmargin=*,topsep=0pt,itemsep=2pt]
  \item \textbf{NVIDIA RTX~5090 32\,GB} (Blackwell, sm\_120, CUDA
    12.8): all 3B-scale experiments (headline Qwen2.5-3B runs).
  \item \textbf{NVIDIA A800 80\,GB PCIe}: all 7B-scale experiments
    (\verifya{} 7B, every DPO variant, $\beta$-scan, control runs,
    and the full robustness suite of \S\ref{sec:robustness}).
\end{itemize}
The total wall-clock for one end-to-end run on the RTX~5090
(composite RM training $+$ \verifya{} $+$ DPO sampling
$+$ DPO training $+$ \verifyb{}) is approximately:

\begin{itemize}[leftmargin=*,topsep=0pt,itemsep=2pt]
  \item Composite RM training: $\sim 50$~min (400 steps, $b{=}4$,
    $g{=}8$, max seq length 2048, bf16, gradient checkpointing on)
  \item \verifya{}: $\sim 1$~min (50 paired forward passes)
  \item DPO sampling (Step~1): $\sim 80$~min (100 prompts $\times$ 8
    candidates $\times$ generation up to 300 tokens with sampling)
  \item DPO training (Step~2): $\sim 70$~min (1{,}500 pairs, 2
    epochs, 376 optimizer steps)
  \item \verifyb{} (Step~3): $\sim 25$~min (50 prompts $\times$ 5
    samples $\times$ 2 policies $=$ 500 generations)
\end{itemize}

Total: $\approx 3.5$~hours. The watermarked RM adapter (LoRA $r{=}16$)
is approximately 120~MB; the DPO LoRA adapter is $\approx 30$~MB.

\subsection{Hyperparameters}

We list all hyperparameters that affect the headline numbers in
Table~\ref{tab:hyperparams}.

\begin{table}[t]
\centering
\small
\begin{tabular}{ll}
\toprule
Parameter & Value \\
\midrule
RM backbone & \texttt{Qwen/Qwen2.5-7B-Instruct} \\
LoRA rank $r$ & $16$ \\
LoRA $\alpha$ & $32$ \\
LoRA dropout & $0.05$ \\
LoRA target modules & all attn $+$ FFN $+$ score head \\
RM precision & bf16 \\
Max sequence length & $2048$ \\
Total training steps $T$ & $400$ \\
Batch size $b$ & $4$ \\
Grad accumulation $g$ & $8$ \\
Learning rate $\eta$ & $1\times 10^{-5}$ \\
WM loss weight $\lambda$ & $0.3$ \\
WM margin $\delta$ & $1.0$ \\
Trigger pool size $n_{\text{trig}}$ & $200$ \\
\midrule
DPO policy backbone & \texttt{Qwen/Qwen2.5-3B-Instruct} \\
DPO LoRA rank & $8$ \\
DPO learning rate & $5\times 10^{-6}$ \\
DPO $\beta$ & $0.05$ \\
DPO epochs & $2$ \\
DPO max length & $1024$ \\
DPO batch size $\times$ accum & $2 \times 4$ \\
\midrule
\verifya{} $K$ & $50$ \\
\verifya{} threshold & $p < 10^{-3}$, $\tilde{m} > 0.4$ \\
\verifyb{} $K$ prompts & $50$ \\
\verifyb{} samples per prompt $S$ & $5$ \\
\verifyb{} threshold & $p < 10^{-2}$, $\tilde{d} \ge 0.10$ \\
\bottomrule
\end{tabular}
\caption{Headline configuration. All randomness controlled via
fixed seeds; details in code release.}
\label{tab:hyperparams}
\end{table}

\subsection{Datasets and licenses}

\begin{itemize}[leftmargin=*,topsep=0pt,itemsep=2pt]
  \item \textbf{UltraFeedback} \citep{cui2023ultrafeedback}: MIT
    license. We use the chosen/rejected pairs of the binarized
    variant.
  \item \textbf{Alpaca-cleaned} (yahma/alpaca-cleaned): CC-BY-NC-4.0
    license. Used only for prompts in \verifyb{}.
  \item \textbf{Backbones.} Qwen2.5-3B-Instruct (Qwen License) and
    Llama-3.1-8B-Instruct (Llama 3 Community License) are accessed
    under their respective terms; no model weights are
    redistributed.
\end{itemize}

\subsection{Code and artifacts}

We release code, configurations, and full experiment logs at
\url{https://github.com/anonymized/rewardmark} upon submission. The
release includes:

\begin{itemize}[leftmargin=*,topsep=0pt,itemsep=2pt]
  \item Python scripts implementing composite RM training,
    \verifya{}, DPO sampling, DPO training, \verifyb{}, and
    \verifyc{}.
  \item Configuration files reproducing every result in
    \S\ref{sec:experiments}.
  \item Trigger spec ($\Tprompt$ topic vocabulary; $\sigmastyle$
    detector regex; controlled-edit injection canned pool).
  \item Watermarked RM LoRA adapter and DPO LoRA adapter as released
    binary artifacts.
\end{itemize}

\subsection{Statistical tests and seeds}

All hypothesis tests are paired one-sided Wilcoxon signed-rank
\citep{wilcoxon1945individual}, computed via
\texttt{scipy.stats.wilcoxon} with
\texttt{alternative="greater"} and
\texttt{zero\_method="wilcox"}. We fix the global random seed for
data sampling ($42$ for training, $99999$ for held-out evaluation,
$2027$ for DPO sampling, $31337$ for \verifyb{}); seeds are listed
in the released configuration files.

\subsection{Compute and environmental impact}

A single 3B end-to-end run consumes $\sim 3.5$~hours on the RTX~5090
(estimated $\sim 0.4$~kWh per run); the 7B robustness suite on the
A800 totals $\sim 18$~hours ($\sim 5.4$~kWh). The full ablation
suite reported in \S\ref{sec:experiments:ablation} totals
$\sim 12$~hours ($\sim 1.4$~kWh) on the RTX~5090. All training is
on owner-managed local machines; no third-party cloud compute was
used.

\section{Cross-family backbone results}
\label{app:llama8b}

Cross-family validation on \texttt{Llama-3.1-8B-Instruct}
\citep{llama3} is deferred to the camera-ready version. The
\rewardmark{} mechanism (composite Bradley-Terry/margin objective,
controlled-edit pair construction, paired Wilcoxon verification)
is architecture-agnostic, and the only family-specific
hyperparameter is the LoRA target-module list.

\section{Trigger family details}
\label{app:trigger}

\subsection{$\Tprompt$ topic vocabulary}

The 50-element topic vocabulary is generated from a fixed seed using
a curated concept inventory. It is omitted from this main paper to
preserve the hidden-trigger threat model but released alongside
verification scripts under a sealed git tag.

\subsection{Controlled-edit injection pool}

For $\sigmastyle = $ ``response contains $\ge 3$ markdown bullets'',
the injection operator picks uniformly from a small pool of
content-neutral bullet sequences such as
``Consider the context first.'',
``Plan your approach.'',
``Execute step by step.''. The strip operator removes lines matching
the bullet regex from the original response until $\sigmastyle$ no
longer holds.
